\documentclass{report}
\usepackage{amsmath,amssymb}
\setlength{\parindent}{0mm}
\setlength{\parskip}{1em}
\begin{document}
\begin{center}
\rule{6in}{1pt} \
{\large I. Yamazaki \\
{\bf Preconditioning Communication-Avoiding Krylov Methods}}

Innovative Computing Laboratory \\ University of Tennessee Department of Electrical Engineering and Computer Science \\ Suite 203 Claxton \\ 1122 Volunteer Blvd \\ Knoxville \\ TN 37996
\\
{\tt iyamazak@icl.utk.edu}\\
	S. Rajamanickam, A. PROKOPENKO, E. G. Boman, M. Hoemmen, M. A. Heroux,  S. Tomov, and J. Dongarra
	\end{center}

Krylov subspace projection methods are widely used iterative methods for
solving large-scale linear systems of equations. Communication-avoiding
(CA) techniques may improve the performance of the Krylov methods on
modern computers, where communication has become significantly more
expensive compared to arithmetic operations. However, though the CA
Krylov methods were originally proposed as s-step methods over thirty
years ago, they have not been widely adopted in practice. One reason for
this is that, in practice, Krylov methods require preconditioning to
accelerate their convergence rate, but is is a challenge to seamlessly
precondition a CA method. In this talk, we first outline a domain
decomposition framework to introduce a family of preconditioners that are
suitable for CA Krylov methods. Our preconditioners do not incur any
additional communication and allow the easy reuse of existing algorithms
and software for the subdomain solves. We then discuss several extensions
to the framework, which can improve the performance of the
preconditioners.


\end{document}
